% New English translation from checked cote 25, site 107 (PDF page 108).
% AI-translated by GPT-6 Astra (OpenAI, ChatGPT Pro mode).
\scanpage{25/107}
\mspage{10}

When D $\hand{\geq} 0$, this implies $n \hand{\geq} 0$, since if \var{\inserted{\unsure{we had|did not have}}}{\unsure{we had|did not have}}{insertion}{ega25_25_107-v0001} $n \hand{\unsure{<|\leq}} 0$ every section of
$\underline{L}$ on X would be zero. By 1.1., if D is a relative divisor
$\hand{\geq} 0$ on P, then there exists a unique decomposition of D as a disjoint
sum of opens $S_n (n \hand{\in} \underline{N})$ such that for every $n \hand{\in} \underline{N}$, $\underline{L}\hand{|}P\hand{|}S_n$ is of
degree n. This reduces the determination of the set of relative divisors
$\hand{\geq} 0$ to the case of relative divisors $\hand{\geq} 0$ of given degree n. (This
\var{\margin{$/\!/$}}{\placedmargin{$/\!/$}}{marginal note placed}{ega25_25_107-v0003} \unsure{text|text} will evidently replace the ``disregarding'' above). This said,
the preceding reflections establish

\uline{Proposition} 2.2. Under the preceding hypotheses, we have a one-to-one
correspondence between the set $\mathrm{Div}^{\unsure{n|m}}(\unsure{P|E|R}/S)$ of relative divisors $\hand{\geq} 0$ of
degree n on P, and the set of invertible quotient Modules $\underline{M}$ of
$\underline{\mathrm{Sym}}^{n}(\underline{E})^{\hand{\vee}}$, (or equivalently, the invertible sub-Modules that are loc
direct summands $\overset{\hand{\vee}}{\underline{M}}$ of $\underline{\mathrm{Sym}}^{n}(\underline{E})$). If D and $\underline{M}$ correspond, then
$\underline{O}_P(D)$ is canonically isomorphic to $\underline{M}\otimes\underline{O}_P(n)$, and the section $s_D$ is
identified by this isomorphism with what one \unsure{guesses|becomes}, taking account of $(*)$.

Note that this description is compatible with every base
change in S. Taking account of the interpretation of inv quotient Modules of
$\underline{\mathrm{Sym}}^{n}(\underline{E})^{\hand{\vee}}$ as the sections over S of $P(\underline{\mathrm{Sym}}^{n}(\underline{E})^{\hand{\vee}})$, we find:
\uline{Corollary} 2.3. The subfunctor \var{\deleted{\var{\xout{\illegible{1}}}{}{typewriter strike-out}{ega25_25_107-v0005}}}{}{deletion}{ega25_25_107-v0004} $\mathrm{Div}^{n}_{P/S}$ of $\mathrm{Div}^{+}_{\unsure{P|E|R}/S}$ is canonically representable
by $P(\underline{\mathrm{Sym}}^{n}(\underline{E})^{\hand{\vee}})$.

Taking account of the considerations of 2.1., it follows:
\uline{Corollary} 2.4. $\underline{\mathrm{Div}}^{+}_{\unsure{P|E|R}/S}$ is canonically representable \var{\xout{between}}{}{typewriter strike-out}{ega25_25_107-v0006} by
the S-prescheme sum of the $P(\underline{\mathrm{Sym}}^{n}_{\underline{P}/S}\unsure{\var{\xout{(}}{}{typewriter strike-out}{ega25_25_107-v0007}|I^{\hand{*}}_{\hand{*}}}\underline{E})^{\hand{\vee}})$, $n \hand{\in} \underline{N}$.
2.5. Suppose now given a finite family $(\underline{E}_i)_{i \hand{\in} I}$ of locally free
Modules on S, whence $P_i = P(\underline{E}_i)$ and a P= fiber product of the
$P_i$, also called the \uline{multiprojective bundle} over S defined by $(\underline{E}_i)$. To
simplify the notation, denote by $\underline{O}_i(n)$ the inverse image on P of the
\tnote{25-107-1}{The source says decomposition of D into opens S-subscript-n, and L on X in the opening argument; these letters are retained. The two spellings tèxte and texte both give text. The alternatives devine and devien are represented by guesses and becomes, leaving the doubtful wording explicit rather than selecting an intended phrase. Inv abbreviates invertible. The canceled xx within relatifs has no English spelling effect.}

\mspage{107}
